When a large language model is shown a novel 2D shape rendered as a text
grid, can it verbally address that shape---name what is at a cell, locate a
token, count occurrences, find the densest row---without producing an
explicit reasoning trace? We design \addrtask, a controlled
$2{\times}2{\times}5$ factorial study over two models (\gpt and \haiku), two
prompting regimes (\Direct and \COT), and five levels of in-context
repetition ($K{=}0,1,2,4,8$ same-grid demonstrations), evaluated on 80
procedurally-generated 7$\times$7 ASCII grids and five question types with
mechanical ground truth.
We find that \Direct prompting plateaus at 79--93\% accuracy even with
$K{=}8$ demos, while \COT lifts both models to 99--100\% from $K{=}0$.
A per-question-type breakdown shows the gap is driven almost entirely by
two scanning queries (\qcount and \qrowmax): random-access queries
(\qpoint, \qlookup, \qneighbor) are at near-ceiling under \Direct prompting.
Paired McNemar tests confirm the \Direct--\COT gap at every $(\text{model},
K)$ cell ($p{\le}0.031$, nine of ten significant under Bonferroni at
$\alpha{=}0.005$). Same-grid repetition matters under \Direct
($0.65 \to 0.93$ for \haiku, $p{<}10^{-6}$), but cannot close the
aggregation gap that \COT closes for free. The picture refines the
``curse of CoT'' story for in-context learning: reasoning helps strictly
when the addressing query requires a scan, and is the cheapest reliability
lever for any LLM workflow grounded in small text grids.
